\subsection{Book の低三角形：混合枝と非零 \texorpdfstring{$D$}{D} 周期への帰着}
\label{nb20:section}

$G_{5,13}=\{01,02,03,04,12,13,14\}$ の二価 cores $2,3,4$ を扱う．
各 core の canonical 非零低低 entry を $-L$ とすると，低三角形の node 総数五から
$L\in\{1,3,4\}$ となる．高 ray $0$ へ正 entry を持つ低 row の先へ
各 core から矢印を引くと，有向三角形になる．
全正枝と全負枝は排除され，存在し得る低三角形は負枝を一個または二個含む．
さらに全低 $D$ は非零で，各 core は一正一負の低 $D$ を持つ．
以下の有限完成法と二つの排除は，高中心 $0,1$ の完成を仮定しない．

\paragraph{有限な node offset．}
有向順序で next row を $(a_i,-b_i,-c_i)$，
previous row を $(-d_i,e_i,-f_i)$ と書くと，
\[
 b_i=a_i+1-c_i,\quad d_i=e_i+1-f_i,\quad
 c_if_i=0,\quad c_i+f_i\in\{1,3,4\}.
\]
三本の低辺の node 条件は
\[
 c_i+f_{i+1}\le4,\quad e_{i+1}=a_i+t_{i+1},\quad
 f_{i+1}-2\le t_{i+1}\le2-c_i
\]
で表示できる．混合枝を調べる際は，高 ray の canonical な交換と
低 $D$ の符号をこの有限 offset に同時に組み込む．
高辺 $01$ があるので，負枝を禁じる低 wedge の補題は適用できない．

\paragraph{全負枝に対する証明された上界と完全排除．}
負枝では $S=T=w$，$H=r-w$ とし，
$w,H\ge2$，$w^2/2<H(H+1)<2w^2$ である．
全負枝の canonical 高 ray の向きが一定の場合は既に排除される．
残る向きを高中心の交換と巡回添字の選択で $(+,+,-)$ にすると，
\[
 w_2=r_1+t_2,\quad r_3=r_2+t_3,\quad w_1=w_3+t_1,
\]
\[
 -1\le t_2\le2,\quad -2\le t_3\le2,\quad -1\le t_1\le1
\]
となる．$17w/10-1/2<r<17w/7$ を二回伝播させると
\[
 \frac{289}{100}w_1-\frac{101}{20}<r_3
 <\frac{17}{7}(w_1+1)
\]
なので $w_1<5235/323<17$，従って
\[
 w_1\le16,\qquad w_3\le17,\qquad w_2\le40
\]
を得る．これは全負枝の仮定から導いた上界であり，混合枝の上界ではない．

固定した $(w,H)$ に対して
\[
 c=\frac{sH}{r},\quad e=\frac{z(H+1)}{r+1},\quad
 E=1-\frac{w^2}{H(H+1)}=\frac MN
\]
と置くと，$c\ge2,e\ge1,E\ne0$，$Ece=e-c+1$ であり，
\[
 (N-Mc)(N+Me)=N(N-M)>0
\]
の両因子は正である．この全正約数から全 scale を復元できるため，
$s,z$ および高分母に探索上限を置く必要はない．
上の reverse 範囲には $91$ 型があり，全 scale を持つ $25$ 型から
$64$ 個の三模型組が生じる．その全てで第一低辺の $D$ 残差が非零となる．
\nolinkurl{verify_book_all_negative_v20.py} と
\nolinkurl{book_all_negative_verification_v20.json} に，
元の row/pair，正の weights，node，$D$ からの全検算を保存した．

\paragraph{低 $D$ の零は完全に除ける．}
一般二高標準形の低係数は，$\delta=T-S$ として
\[
 D_C=\delta-\frac Sr s+\frac T{r+1}z,\qquad D_D=s-z-1
\]
である．二つの低 $D$ は同時に負にならない．
正枝では $r\delta>S$ であり，
$D_D<0$ なら $z\ge s$ から $D_C>0$，
$D_D=0$ なら
\[
 D_C=\frac{(r+z+1)(r\delta-S)}{r(r+1)}>0
\]
となる．負枝の $E<0$ では $D_C<0<D_D$，
$E>0$ では $D_D<0\le D_C$ である．後者の等号には
$e-c=1,c=H$ が必要で，$w^2=H(H+1)-2$ から
\[
 (r,s,z,S,T)=(4,4,5,2,2),\qquad(D_C,D_D)=(0,-2)
\]
だけが残る．従って二低 $D$ 同時零はなく，負と零の組はこの一模型だけである．
この局所零分類には分母の奇数性は不要である．

この唯一の模型を core $2$，canonical order を $(0,1,3,4)$ とすると，
$q=(9,11,1,1)$ であり，低辺 $23,24$ の完全な相手集合はそれぞれ四個，一個である．
core $4$ は $(r,s,z,S,T)=(2,2,1,3,6)$，order $(0,1,2,3)$ に固定される．
core $3$ と最後の辺 $34$ は次表のとおりである．node の順は $(0,1,2)$ とする．
\begin{center}
\small
\begin{tabular}{lccc}
core $3$ の parameters&canonical order&node 三つ&$D^3_4+D^4_3$\\\hline
$(2,2,1,3,6)$&$(0,1,4,2)$&$(2,-2,5)$&$2$\\
$(3,3,2,4,6)$&$(0,1,4,2)$&$(1,-1,5)$&$1$\\
$(4,4,5,2,2)$&$(1,0,2,4)$&$(2,-3,6)$&$-2$\\
$(4,4,2,7,10)$&$(1,0,2,4)$&$(-6,2,9)$&$1$
\end{tabular}
\end{center}
四組とも完成しない．復元は全二十一 node 行と指定低 $D$ から行い，
scale や分母の数値上限を使わない．
\nolinkurl{verify_book_low_D_cycle_v20.py} はこの表を独立に再計算する．

二本以上の零低辺があれば，交点 core の二低 $D$ がともに零になって矛盾する．
零低辺が一本なら，その両端の残る低 $D$ がともに正である場合，
第三 core の二低 $D$ がともに負となる．それ以外なら負と零の唯一の core が現れ，
上の四組の排除に反する．従って全低 $D$ は非零である．
三辺の相殺は正端三つ・負端三つを与え，一つの core に負端二つはないため，
各 core は一正一負となる．各辺を負端から正端へ向けると有向三角形になる．
この向きと高 ray $0$ への正 entry の向きは異なる定義であり，
同方向・逆方向の両場合を今後保持する．

\paragraph{固定 core の有限完全完成法．}
一つの二価 core の行列を固定すると，各低辺の二十一 node 分配は相手の低 row 全体を固定する．
対向 $D$ はその低 ray の $D=d$ も固定する．canonical reverse data $r,S,T$ は，
零低低 row なら $r$ の固定と $L\le4$，非零低低 row なら $S,T$ の固定と
負枝の $w^2/2<H(H+1)<2w^2$ により有限個となる．

$h=S-r$，$\delta=T-S\in\{0,2,3\}$ とすると，残る scale は
\[
 Fsz+As+Bz+C=0,\qquad \ell s+mz+n=0,
\]
\[
 F=(2r+1)h-r^2(\delta-1),\quad
 A=h(r+1),\quad B=-r(h+\delta-1),\quad C=r(r+1)
\]
を満たす．低 $D$ の条件は
\[
 (\ell,m,n)=
 \begin{cases}
 (-S(r+1),rT,(\delta-d)r(r+1)),&C\text{ ray},\\
 (1,-1,-d-1),&D\text{ ray}.
 \end{cases}
\]
$\gcd(2r+1,r^2)=1$ と $\delta\in\{0,2,3\}$ から $F\ne0$，
また $\ell m\ne0$ なので，消去した二次式
\[
 -F\ell s^2+(-Fn+Am-B\ell)s+(Cm-Bn)=0
\]
の最高次係数は非零である．固定 reverse data ごとに高々二つの scale が残り，
整数平方判定と整除から全て復元できる．七・十三の共鳴もこの二次式に含まれる．
二つの低隣接 core の有限集合を作り，残る低辺を検査すれば全低三角形完成を得る．
固定 core 自体が動く無限集合を処理したことにはならない．

実装は \nolinkurl{research/book_low_triangle_completion_v19.py} であり，
独立 signed-divisor catalog との有界な実装照合は
\nolinkurl{verify_book_low_triangle_v19.py} に保持する．
再探索の重複を避ける資料として，負枝 $2\le w\le150$ の $1259$ canonical cores には
低三角形完成がなかったことを残す．両低 ray に相手を持った六模型は
\nolinkurl{research/book_negative_w150_candidates_v19.json} にある．
これは混合枝の $w>150$ を排除しない．次はこの上限を増やすことより，
混合枝・有限 offset・非零 $D$ 周期の三条件を組み合わせた代数的比較を優先する．
